% Sources: Neso videos (monoalphabetic; Playfair parts 1-2; Playfair solved question; crypto/notes3.md), college Lecture 5
% (Forouzan 3.2: monoalphabetic key table, statistical attack, Playfair secret key and "hello", cryptanalysis of Playfair),
% Stallings 2.2 (frequency analysis, homophones, Playfair strength). Every number recomputed by script (scratchpad/r2/verify.py).
\providecommand{\NESOACADEMY}{{N,E,S,O,A},{C,D,M,Y,B},{F,G,H,{I/J},K},{L,P,Q,R,T},{U,V,W,X,Z}}
\providecommand{\JUPITER}{{{I/J},U,P,T,E},{R,A,B,C,D},{F,G,H,K,L},{M,N,O,Q,S},{V,W,X,Y,Z}}
\providecommand{\NESOAPP}{{N,E,S,O,A},{P,B,C,D,F},{G,H,{I/J},K,L},{M,Q,R,T,U},{V,W,X,Y,Z}}
\providecommand{\PFFOROUZAN}{{L,G,D,B,A},{Q,M,H,E,C},{U,R,N,{I/J},F},{X,V,S,O,K},{Z,Y,W,T,P}}

\section{Monoalphabetic Cipher}\label{sec:monoalpha}

The Caesar (shift) cipher has two weaknesses that make it useless. Its key space has only 25 useful keys, so an attacker simply tries them all. Worse, every letter is shifted by the \emph{same} amount, so the shift found for a single letter decrypts the whole message. The monoalphabetic cipher removes both weaknesses at once: instead of shifting the alphabet, it replaces each letter by an \emph{arbitrary} other letter, chosen independently for each of the 26 letters. The only constraint is that the replacement must be one-to-one, so that the receiver can undo it.

\subsection{Permutations}
A \imp{permutation of a finite set $S$ is an ordered sequence of all the elements of $S$, with each element appearing exactly once.} The set $S = \{a, b, c\}$ has six permutations: \ttx{abc, acb, bac, bca, cab, cba}. In general a set of $n$ elements has $n! = n(n-1)\cdots 2 \cdot 1$ permutations, because the first position can be filled in $n$ ways, the second in $n - 1$ ways (one element is used up), and so on down to one way for the last position. For example $2! = 2$ and $4! = 24$.

Permutations matter here because the key of a monoalphabetic cipher \emph{is} a permutation of the alphabet: writing the ciphertext letters for $a, b, \dots, z$ in order gives a sequence in which each of the 26 letters appears exactly once. The number of possible keys is therefore the number of permutations of 26 letters, $26!$.

\subsection{The cipher}
In a \textbf{monoalphabetic substitution cipher} Alice and Bob agree on a table that pairs every plaintext letter with one ciphertext letter. The bottom row of the table, the \textbf{cipher line}, can be any permutation of the 26 letters. To encrypt, each plaintext letter is looked up in the top row and replaced by the letter below it; to decrypt, each ciphertext letter is found in the bottom row and replaced by the letter above it. The relationship between plaintext and ciphertext symbols is \textbf{one-to-one}: a plaintext letter always becomes the same ciphertext letter, and no two plaintext letters share a ciphertext letter. If two letters shared one, the receiver could not tell which was meant, so the table would not be decryptable.

The key space has $26! \approx 4.03 \times 10^{26} \approx 2^{88}$ keys, compared with 25 for the shift cipher. A computer testing $10^{12}$ keys per second would need about 13 million years to try them all, so the slide's conclusion is that the large key space ``would seem to eliminate brute-force techniques''. The shift ciphers are just 26 of these $26!$ permutations, those in which the cipher line is the alphabet rotated. Although a permutation is involved, the monoalphabetic cipher is still a \textbf{substitution} cipher: the permutation is applied to the \emph{alphabet} to build the key, while the letters of the message keep their positions. Rearranging the message itself would be a transposition.

\textbf{Example (lecture).} Under one key, ``Neso'' becomes \ttx{DULA}. The distance each letter moves is different, so learning one pair of the table says nothing about the others; with Caesar the distance is always 3.
\begin{center}\small\setlength{\tabcolsep}{6pt}
\begin{tabular}{@{}lcccc@{}}
\toprule
Plaintext & N (13) & E (4) & S (18) & O (14)\\
Ciphertext & \textbf{D} (3) & \textbf{U} (20) & \textbf{L} (11) & \textbf{A} (0)\\
Shift $(C-p) \bmod 26$ & 16 & 16 & 19 & 12\\
\bottomrule
\end{tabular}
\end{center}

\begin{worked}[: the slide's key table (Forouzan)]
\textbf{Q.} Using the key below, encrypt ``this message is easy to encrypt but hard to find the key'', and show how Bob decrypts.
\begin{center}\small\setlength{\tabcolsep}{2.4pt}
\begin{tabular}{@{}l*{26}{c}@{}}
\toprule
Plain & a&b&c&d&e&f&g&h&i&j&k&l&m&n&o&p&q&r&s&t&u&v&w&x&y&z\\
Cipher & N&O&A&T&R&B&E&C&F&U&X&D&Q&G&Y&L&K&H&V&I&J&M&P&Z&S&W\\
\bottomrule
\end{tabular}
\end{center}
\textbf{A.} Remove the spaces and replace each letter by the one below it: t$\to$I, h$\to$C, i$\to$F, s$\to$V, so ``this'' $\to$ \ttx{ICFV}; ``message'' $\to$ \ttx{QRVVNER}; and so on. The ciphertext is
\begin{center}\small\ttfamily ICFVQRVVNERFVRNVSIYRGAHSLIOJICNHTIYBFGTICRXRS\end{center}
Bob finds each ciphertext letter in the bottom row and reads the letter above it (I$\to$t, C$\to$h, F$\to$i, V$\to$s, \dots), which recovers the plaintext. Note that the three \ttx{s} of ``message is'' all became \ttx{V} and every \ttx{e} became \ttx{R}: the frequency pattern of the plaintext survives, which is the subject of the next subsection.
\end{worked}

\subsection{Breaking it: frequency analysis}
The huge key space does not protect the cipher, because the attacker never searches through keys. A monoalphabetic cipher only \emph{renames} letters. If the key maps $x$ to $\pi(x)$, then every occurrence of $x$ in the plaintext becomes an occurrence of $\pi(x)$ in the ciphertext and nothing else becomes $\pi(x)$, so $\pi(x)$ occurs in the ciphertext exactly as often as $x$ occurs in the plaintext. The histogram of ciphertext letters is the histogram of plaintext letters with the bars relabelled. \imp{A monoalphabetic cipher does not change the frequency of characters}, so each letter of the key can be found separately by matching frequencies, and the attack needs a few dozen decisions instead of $26!$ guesses. This is a \textbf{statistical} (frequency) attack, and it needs only the ciphertext, which makes it a ciphertext-only attack, the most likely kind in practice.

The attack works because English letters are far from equally common. The relative frequencies (lecture chart, from a large corpus) are:
\begin{center}
\begin{tikzpicture}[x=6.1mm,y=2.5mm]
  \foreach \L/\v [count=\i] in {A/8.167,B/1.492,C/2.782,D/4.253,E/12.702,F/2.228,G/2.015,H/6.094,I/6.996,J/0.153,K/0.772,L/4.025,M/2.406,N/6.749,O/7.507,P/1.929,Q/0.095,R/5.987,S/6.327,T/9.056,U/2.758,V/0.978,W/2.360,X/0.150,Y/1.974,Z/0.074}{
    \fill[acc!55] (\i-0.33,0) rectangle (\i+0.33,\v);
    \node[font=\scriptsize,below] at (\i,0) {\L};
  }
  \fill[warn] (5-0.33,0) rectangle (5.33,12.702) node[above,font=\scriptsize,xshift=-1.1mm] {12.7};
  \fill[warn] (20-0.33,0) rectangle (20.33,9.056) node[above,font=\scriptsize,xshift=-1.1mm] {9.1};
  \draw (0.4,0) -- (26.6,0);
  \draw (0.4,0) -- (0.4,13.5);
  \foreach \y in {0,4,8,12} \node[font=\scriptsize,left] at (0.4,\y) {\y};
  \node[font=\scriptsize,rotate=90] at (-1.1,6.5) {\% of letters};
\end{tikzpicture}
\end{center}
Ranked from most to least common, the letters are
\begin{center}
\ttx{e t a o i n s h r d l c u m w f g y p b v k j x q z}.
\end{center}
The letter \ttx{e} alone is 12.7\% of English text and \ttx{t} is 9.1\%; the nine letters \ttx{e t a o i n s h r} together make up about 70\% of it, while \ttx{j, x, q, z} are each below 0.2\%. Pairs and triples are just as uneven. The most common \textbf{digrams} (two-letter groups) are \ttx{th, he, in, er, an, re, ed, on, es, st, en, at, to, nt, ha, nd, ou, ea, ng, as, or, ti, is, et, it, ar, te, se, hi, of}, and the most common \textbf{trigrams} are \ttx{the, ing, and, her, ere, ent, tha, nth, was, eth, for, dth} (lecture table, from Forouzan).

The attack therefore proceeds as follows. Eve counts the ciphertext letters and guesses that the most frequent one stands for \ttx{e} and the next for \ttx{t}. She then tests the guesses against digrams and trigrams: if \ttx{t} has been placed, a frequent pattern ``\ttx{t?e}'' is almost certainly \ttx{the}, which reveals \ttx{h}. Each new letter makes more words partly readable, and the remaining letters are filled in by recognising words. The statistics only become reliable with enough text: with about 100 letters the top two or three ranks usually match English, while the lower ranks are shuffled and have to be fixed by word guessing. For a very short message the counts may mislead completely.

\begin{worked}[: the lecture's frequency attack (Stallings' ciphertext, 120 letters)]
\textbf{Q.} The following ciphertext was produced by a monoalphabetic cipher. Use a statistical attack to find the plaintext.
\begin{center}\small\ttfamily
UZQSOVUOHXMOPVGPOZPEVSGZWSZOPFPESXUDBMETSXAIZ\\
VUEPHZHMDZSHZOWSFPAPPDTSVPQUZWYMXUZUHSX\\
EPYEPOPDZSZUFPOMBZWPFUPZHMDJUDTMOHMQ
\end{center}
\textbf{A. Step 1: count.} The 120 letters occur as follows (C, K, L, N, R do not occur):
\begin{center}\small\setlength{\tabcolsep}{3pt}
\begin{tabular}{@{}l*{21}{c}@{}}
\toprule
Letter & \textbf{P}&\textbf{Z}&S&U&O&M&H&D&E&V&X&F&W&Q&T&A&B&G&Y&I&J\\
Count & 16&14&10&10&9&8&7&6&6&5&5&4&4&3&3&2&2&2&2&1&1\\
\bottomrule
\end{tabular}
\end{center}
\textbf{Step 2: the two peaks.} \ttx{P} is $16/120 = 13.33\%$ and \ttx{Z} is $14/120 = 11.67\%$, close to the 12.7\% of \ttx{e} and 9.1\% of \ttx{t}. Guess \ttx{P} $=$ e and \ttx{Z} $=$ t.

\textbf{Step 3: digrams and trigrams.} The digram \ttx{ZW} occurs 3 times, and \ttx{ZWP} once. With the guesses this is ``t?e'', so \ttx{ZWP} $=$ ``the'' and \ttx{W} $=$ h. Then \ttx{ZWS} $=$ ``th?'', which fits ``that'', so \ttx{S} $=$ a.

\textbf{Step 4: fill in.} With \ttx{e, t, h, a} placed, $16 + 14 + 4 + 10 = 44$ of the 120 letters (37\%) are known and words become guessable. For example, \ttx{ZWSZ OPFPESX} now reads ``that ?e?e?a?'', which suggests ``that several'' and gives \ttx{O} $=$ s, \ttx{F} $=$ v, \ttx{E} $=$ r, \ttx{X} $=$ l. Each guess is checked against the rest of the text, and the remaining letters follow in the same way. The result is
\begin{center}\small
\emph{it was disclosed yesterday that several informal but direct contacts have been made with political representatives of the viet cong in moscow}
\end{center}
and the recovered key (cipher $\to$ plain) is
\begin{center}\small\setlength{\tabcolsep}{2.6pt}
\begin{tabular}{@{}l*{21}{c}@{}}
\toprule
Cipher & A&B&D&E&F&G&H&I&J&M&O&\textbf{P}&Q&S&T&U&V&\textbf{W}&X&Y&\textbf{Z}\\
Plain  & b&f&n&r&v&y&c&u&g&o&s&\textbf{e}&w&a&m&i&d&\textbf{h}&l&p&\textbf{t}\\
\bottomrule
\end{tabular}
\end{center}
\textbf{Check:} applying this table to all 120 letters gives the sentence above. Note that \ttx{S} (a) and \ttx{U} (i) tie at 10, ahead of \ttx{O} (s) at 9: below the top two ranks, a short text does not follow the English ranking exactly, which is why steps 3--4 are needed.
\end{worked}

\subsection{Pros, cons, countermeasure}
\textbf{Advantages.} The monoalphabetic cipher is as simple to use as the Caesar cipher, but its key space of $26!$ makes brute force impossible, and knowing one letter's substitute reveals nothing about the others. It is therefore much more secure than the Caesar cipher against an attacker who only tries keys.

\textbf{Disadvantages.} Since the ciphertext reflects the frequency data of the original alphabet, the cipher is easy to break with frequency analysis; the guessing attack above needs only a short ciphertext and some knowledge of English letter frequencies. The key is also long (a full table of 26 letters) and awkward to share and memorise, which is why practical versions derive it from a key phrase (see the homework below).

\textbf{Countermeasure: homophones.} A \textbf{homophonic substitution} gives each plaintext letter several possible ciphertext symbols, called \textbf{homophones}, and the encryptor picks one of them at random for each occurrence. The number of homophones is made proportional to the letter's frequency: \ttx{e} (about 13\%) gets about 13 symbols, while \ttx{z} gets one. Since each of the 13 symbols for \ttx{e} is then used for roughly $12.7/13 \approx 1\%$ of the text, every symbol appears with about the same frequency and the peaks in the histogram disappear. The ciphertext alphabet must be larger than 26, for example the numbers 00 to 99. Decryption is still unique, because each symbol belongs to only one plaintext letter.

Homophones flatten the \emph{single-letter} frequencies, but they do not make the cipher secure. Each plaintext letter still affects only one ciphertext symbol, so multi-letter patterns survive: \ttx{t} is still very often followed by \ttx{h}, so the symbols of \ttx{t} are often followed by symbols of \ttx{h}, and the attacker can group the symbols that behave alike and then proceed as before. Two better remedies were developed from this observation: encrypting \emph{several letters together} (Playfair, Hill), and using \emph{several cipher alphabets} chosen by a key (the polyalphabetic ciphers of \S\ref{sec:poly}).

\subsection{Homework: encrypt with a key phrase}
Since a random 26-letter table is hard to remember, the key is often derived from a key phrase: write the phrase without spaces, delete every letter that has already appeared, and use the result as the cipher line. If the phrase does not contain all 26 letters, the unused letters are appended in alphabetical order (exactly as in the Playfair square below).

\begin{worked}[: lecture homework]
\textbf{Q.} Encrypt ``Attack postponed to tomorrow and do not use our secret paper until further info'' with a monoalphabetic cipher whose key is ``The quick brown fox jumps over the lazy dog'' (ignore the second and later occurrences of a letter in the key).

\textbf{A.} Removing spaces and repeated letters from the key phrase leaves \ttx{THEQUICKBROWNFXJMPSVLAZYDG}. The phrase is a pangram, so this already contains all 26 letters and is the whole cipher line:
\begin{center}\small\setlength{\tabcolsep}{2.4pt}
\begin{tabular}{@{}l*{26}{c}@{}}
\toprule
Plain & a&b&c&d&e&f&g&h&i&j&k&l&m&n&o&p&q&r&s&t&u&v&w&x&y&z\\
Cipher & T&H&E&Q&U&I&C&K&B&R&O&W&N&F&X&J&M&P&S&V&L&A&Z&Y&D&G\\
\bottomrule
\end{tabular}
\end{center}
Substituting letter by letter (spaces kept only for readability):
\begin{center}\small\ttfamily
TVVTEO JXSVJXFUQ VX VXNXPPXZ TFQ QX FXV LSU XLP SUEPUV JTJUP LFVBW ILPVKUP BFIX
\end{center}
\textbf{Check} on the first word: a$\to$T, t$\to$V, t$\to$V, a$\to$T, c$\to$E, k$\to$O gives \ttx{TVVTEO}. The doubled \ttx{tt} became the doubled \ttx{VV}, another pattern a monoalphabetic cipher cannot hide.
\end{worked}

%%PLAYFAIR%%
